% Resultados del ejemplo A: fuerzas axiales en kN (tracción en azul, compresión
% en naranja) y reacciones (capítulo 10). Los valores vienen de recursos/valores.tex.
\begin{tikzpicture}[scale=0.8,
    reaccion/.style={-{Stealth[length=7pt]}, line width=1.2pt, color=principal},
    valor/.style={font=\sffamily\small, fill=white, inner sep=1.5pt,
      rounded corners=1pt}]
  \coordinate (n1) at (0,0);
  \coordinate (n2) at (6,0);
  \coordinate (n3) at (3,4);
  \draw[barra traccion, line width=2.4pt] (n1) -- node[valor, text=traccion]
    {\valor{c10_n1}} (n2);
  \draw[barra compresion, line width=2.4pt] (n2) -- node[valor, text=compresion]
    {\valor{c10_n2}} (n3);
  \draw[barra compresion, line width=2.4pt] (n1) -- node[valor, text=compresion]
    {\valor{c10_n3}} (n3);
  \apoyoarticulado{(n1)}
  \apoyorodillo{(n2)}
  \cargapuntual{(n3)}{(1.4,0)}{\SI{20}{\kilo\newton}}
  \cargapuntual{(n3)}{(0,-1.4)}{\SI{50}{\kilo\newton}}
  % Reacciones
  \draw[reaccion] (-0.45,-0.25) -- ++(-1.3,0)
    node[left, font=\small] {\valor{c10_r1_magnitud}};
  \draw[reaccion] (-0.45,-1.7) -- ++(0,1.15)
    node[pos=0, below, font=\small] {\valor{c10_r2_magnitud}};
  \draw[reaccion] (6,-1.9) -- ++(0,1.15)
    node[pos=0, below, font=\small] {\valor{c10_r4_magnitud}};
  \foreach \n in {n1, n2, n3}
    \node[nodo] at (\n) {};
  \node[etiqueta nodo, above left=2pt] at (n1) {1};
  \node[etiqueta nodo, above right=2pt] at (n2) {2};
  \node[etiqueta nodo, above right=4pt] at (n3) {3};
  % Leyenda
  \draw[barra traccion, line width=2.4pt] (7.6,3.6) -- ++(0.8,0)
    node[right, font=\small, color=black] {tracción};
  \draw[barra compresion, line width=2.4pt] (7.6,3.0) -- ++(0.8,0)
    node[right, font=\small, color=black] {compresión};
  \draw[reaccion] (7.6,2.4) -- ++(0.8,0)
    node[right, font=\small, color=black] {reacción};
  \node[font=\small, anchor=west] at (7.6,1.8) {valores en kN};
\end{tikzpicture}
